\documentclass[12pt]{article}

\usepackage[dvips,letterpaper,margin=0.75in,bottom=0.5in]{geometry}
\usepackage{cite}
\usepackage{slashed}
\usepackage{graphicx}
\usepackage{amsmath}
\usepackage{amssymb}
\usepackage{braket}
\usepackage{pythonhighlight}



\begin{document}
\newcommand{\ihbar}{\ensuremath{i \hbar}}
\newcommand{\Pss}{\ensuremath{\Psi^*}}
\newcommand{\dPsidt}{\ensuremath{ \frac{\partial \Psi}{\partial t} }}
\newcommand{\dPsidx}{\ensuremath{ \frac{\partial \Psi}{\partial x} }}
\newcommand{\ddPsidx}{\ensuremath{ \frac{\partial^2 \Psi}{\partial x^2} }}
\newcommand{\dPssdt}{\ensuremath{ \frac{\partial \Psi^*}{\partial t} }}
\newcommand{\dPssdx}{\ensuremath{ \frac{\partial \Psi^*}{\partial x} }}
\newcommand{\ddPssdx}{\ensuremath{ \frac{\partial^2 \Psi^*}{\partial x^2} }}

\newcommand{\dphidt}{\ensuremath{ \frac{d \phi}{dt} }}
\newcommand{\dpsidx}{\ensuremath{ \frac{d \psi}{dx} }}
\newcommand{\ddpsidx}{\ensuremath{ \frac{d^2 \psi}{dx^2} }}

\title{PHY 115L \\ The Gaussian Wave Packet}
\author{Michael Mulhearn}

\maketitle

\section{Introduction}

For this assignment, we will animate solutions to the 1-D
time-dependent Schr\"odinger Equation (TDSE):
\begin{equation}
\label{eqn:se}
\ihbar \frac{\partial \Psi}{\partial t} = H \Psi(x,t)
\end{equation}\\[3pt]
where $H$ is the Hamiltonian operator:
$$H = - \frac{\hbar^2}{2 m} \frac{\partial^2}{\partial x} + V(x) $$
The energy eigenstates of the Hamiltonian:
$$H \ket{n} = E_n \ket{n}$$
are a basis for the Hilbert space and have a well understood time-dependence:
$$\ket{n(t)} = \exp{\left(- \frac{i E_n t}{\hbar}\right)} \ket{n}$$
So once we know the coefficients of the energy eigenstates for the wave function at $t=0$:
$$c_n = \braket{n|\Psi(0)}$$
then we have the time-dependent solution as well:
\begin{equation}
  \label{eqn:expansion}
 \ket{\Psi(t)} = \sum_n c_n \exp{\left(- \frac{i E_n t}{\hbar}\right)} \ket{n}
\end{equation}

\section{System of Units}

As ever, we will employ a numerical system with a characteristic
length scale $a$.  In numerical analysis, it is more natural to use
the wave number $k$, related to the momentum $p$, by:
$$p = \hbar k$$
In our system of units $k$ is measured in units of $1/a$.
Equivalently, we can consider momentum to be measured in units of:
$$p_0 = \hbar /a $$
Energy will be measured in units of:
\begin{equation}
E_0 \equiv \frac{\hbar^2}{2ma^2}
\end{equation}
As always, quantities like $\hbar$, $a$, $p_0$ and $E_0$ should not
appear in your program.  We keep track of them explicitly for analytic
calculations, where dimensional analysis is useful, but we set them to
one in our computer programs.

\section{Time-Evolution of a Free Particle Wave Packet}

For a free particle, the energy eigenstates are the momentum
eigenstates, which are continuous.  Therefore, the sum in
Equation~\ref{eqn:expansion} is replaced by an integral:
\begin{equation}
\Psi(x,t) = \frac{1}{\sqrt{2\pi}} \int_{-\infty}^{\infty} \; \phi(k) \; \exp\left[ik\left(x-\frac{\hbar k}{2m}t\right)\right]\; dk 
\end{equation}\\[3pt]
where the discrete coefficients have been replaced by the continuous Fourier
transform of the initial state:
\begin{equation}
\phi(k) = \frac{1}{\sqrt{2\pi}} \int_{-\infty}^{\infty} \; \Psi(x,0) \; e^{-ikx} \; dx 
\end{equation}
We will consider the initial state to be:
\begin{equation}
\label{eqn:initial}
\Psi(x,0) = N \exp\left( -\frac{x^2}{4\,\alpha^2} \right) \exp(i\beta x)
\end{equation}
with normalization factor:
$$N = \left( \frac{1}{2 \pi \alpha^2} \right)^\frac{1}{4}$$
Notice that:
$$|\Psi(x,0)|^2 = \frac{1}{\sqrt{2\pi} \alpha} \exp\left( -\frac{x^2}{2\,\alpha^2} \right)$$
which is a Gaussian distribution with width $\alpha$.  The Fourier transform of this initial state is:
\begin{eqnarray*}
\phi(k) &=& \frac{1}{\sqrt{2\pi}} \int_{-\infty}^{+\infty} \; e^{-ikx} \; \Psi(x,0) \, dx \\
        &=& \frac{N}{\sqrt{2\pi}} \int_{-\infty}^{+\infty} \; \exp(-X) \; dx \\
\end{eqnarray*}
where:
\begin{eqnarray*}
  X &=& \frac{1}{4} \left( \frac{x}{\alpha} \right)^2 + i\,(k-\beta)\,x\\
    &=& \frac{1}{4} \left[ \frac{x}{\alpha}  + i 2 \alpha (k-\beta) \right]^2 + \alpha^2(k-\beta)^2\\
\end{eqnarray*}
and we have ``completed the square''.  To compute the integral, we make the substitution:
$$u = \frac{1}{2} \left[ \frac{x}{\alpha} + i 2 \alpha (k-\beta)\right]$$
so:
$$X = u^2 + \alpha^2 (k-\beta)^2$$
and
$$dx = 2 \alpha \, du.$$
Our new integration variable $u$ is complex, and our integral is now a
contour integral on the complex plane, but we will get the correct
result by simply rotating back to the $x$-axis, e.g.:
$$u = \pm\infty + i 2 \alpha (k-\beta) \to \pm\infty$$
so the integral becomes:
$$\phi(k) = M \, \exp(-\alpha^2(k-\beta)^2)$$
where
$$M = \frac{2 \alpha N}{\sqrt{2\pi}} \int_{-\infty}^{\infty} e^{-u^2} du = \left( \frac{2\alpha^2}{\pi} \right)^\frac{1}{4}$$
Notice that:
$$|\phi(k)|^2 = \frac{1}{\sqrt{2\pi} \left(\frac{1}{2 \alpha}\right)} \exp\left( -\frac{(k-\beta)^2}{2\left( \frac{1}{2 \alpha}\right)^2}\right)$$
which is a Gaussian distribution with mean $\beta$ and
$$\sigma_k = \frac{1}{2 \alpha}$$
Recalling that $\sigma_x = \alpha$ and $p =\hbar k$ we have:
$$\sigma_x \sigma_p = \frac{\hbar}{2}$$
so the wave packet has minimum uncertainty.

Next we calculate the time-dependent wave function from:
\begin{eqnarray*}
  \Psi(x,t) &=& \frac{1}{\sqrt{2\pi}}\int_{-\infty}^{+\infty} \phi(k) \; e^{ikx} \; \exp\left(\frac{-iEt}{\hbar}\right) dk \\
  &=& M \int_{-\infty}^{+\infty} \; \exp(-X) \, dk\\  
\end{eqnarray*}
where:
\begin{eqnarray*}
  X &=& \alpha^2(k-\beta)^2 - ikx + i \frac{\hbar t}{2 m}\,k^2
\end{eqnarray*}
and we have used:
$$E = \frac{\hbar^2 k^2}{2 m}.$$
The algebra that follows turns out much simpler if we first make the substitution:
$$q = k - \beta \implies k = q + \beta$$
So we have:
\begin{eqnarray*}
  X &=& \alpha^2 q^2 - i(q+\beta)x + i \frac{\hbar t}{2 m}\,(q+\beta)^2\\
  &=& \left(\alpha^2 + i \frac{\hbar t}{2 \, m}\right) q^2 + i(\frac{\hbar \beta t}{m}-x)q - i\beta x + i\frac{\hbar \beta^2}{2 m}\,t\\
  &=& A \left( \frac{q}{\beta} \right)^2 + B \frac{q}{\beta} - i\beta x + i \frac{\hbar \beta^2}{2 m}\,t\\  
\end{eqnarray*}
where:
\begin{eqnarray*}
  A &=& n^2 + i \omega t \\
  B &=& i (2 \omega t - \beta x) \\
  n &=& \alpha \beta \\
  \omega &=& \frac{\hbar \beta^2}{2m}\\ 
\end{eqnarray*}
Notice the coefficient of $q$ is pure imaginary, which is the
simplification that resulted from our substitution.  Completing the square we find:
\begin{equation*}
X = A \left[ \frac{q}{\beta} + \frac{B}{2A} \right]^2 - \frac{B^2}{4A}-i\beta x+i \omega t
\end{equation*}
and substituting
$$u = \sqrt{A} \left[\frac{q}{\beta} + \frac{B}{2A}\right]$$
we have:
\begin{equation*}
X = u^2 - \frac{B^2}{4A}-i\beta x+i \omega t.
\end{equation*}
and:
\begin{eqnarray*}
-\frac{B^2}{4A} &=& \frac{(2 \omega t - \beta x)^2 }{4 (n^2+i\omega t)} \\[5pt]
&=& \frac{\beta^2}{4 n^2} \, \frac{(x - \frac{2 \omega}{\beta} t)^2}{1 + i (\omega t / n^2)}\\
&=& \frac{(x - v t)^2}{4 \alpha^2 (1 + i (\omega t / n^2))}\\
\end{eqnarray*}
where:
$$v = \frac{2\omega}{\beta} = \frac{\hbar \beta}{m}$$
We also have:
$$-i\beta x + i\omega t = i \beta (x - v t / 2)$$
Noting that:
$$dk = dq = \frac{\beta}{\sqrt{A}} du$$
we have:
\begin{equation}
\label{eqn:result}
\Psi(x,t) = N(t) \, \exp\left( -\frac{(x - v t)^2}{4 \alpha^2 (1 + i (\omega t / n^2))} \right)
\exp\left(i\beta\left(x-\frac{vt}{2}\right)\right)
\end{equation}
where:
\begin{eqnarray*}
  N(t) &=& \frac{1}{\sqrt{2\pi}} \left( \frac{2\alpha^2}{\pi} \right)^\frac{1}{4} \frac{\beta}{\sqrt{A}} \int_{-\infty}^{\infty} e^{-u^2} du\\ 
  &=& \left( \frac{1}{2\pi \, \alpha^2} \right)^\frac{1}{4} \, \frac{1}{\sqrt{1 + i \omega t/n^2}}\\
\end{eqnarray*}
There is a lot of phenomenology contained in Equation~\ref{eqn:result}.  Things you can notice:
\begin{itemize}
\item The packet starts with initial width $\alpha$ and becomes wider over time via dispersion.
\item The packet travels with group velocity $v=\hbar \beta/m$, i.e. the classical velocity of a particle with momentum $\hbar \beta$.
\item The carrier waves (complex exponential) have a phase velocity that is one-half the group velocity.
\end{itemize}
You should arrange your animation so that these features are evident.

\section{Animation in Scientific Python}

To animate the time dependent solutions, we will use the FuncAnimation
feature from matplotlib.animation.  To run in a jupyter notebook, you
will want to use the matplotlib widget backend for animation:
\begin{python}
import numpy as np
import matplotlib.pyplot as plt
from matplotlib.animation import FuncAnimation
%matplotlib widget
\end{python}
Remember that the last line is not a comment, it is ``line magic''
that starts the backend.  If you haven't used the matplotlib widget
before, you may have to install it first.  For example, under conda:
\begin{verbatim}
conda install ipympl
\end{verbatim}
An example animation is provided on the course website (uncertainty.ipynb). In that example, the animation is implemented like this:
\begin{python}
# animation parameters
tstep = 1
fps = 20
frame_ms = int(1000 / fps)
second_tau = fps * tstep

# setup the figures
fig, axes = plt.subplots(1, 2, figsize=(16, 9), dpi=1920/32)
xf = np.linspace(-10, 10, 1000)
la, = axes[0].plot([], [], "b-")
lb, = axes[1].plot([], [], "r-")

axes[0].set_xlim(-10, 10)
axes[1].set_xlim(-10, 10)
axes[0].set_ylim(0, 1.5)
axes[1].set_ylim(0, 1.5)

axes[0].set_xlabel("$x$")
axes[0].set_ylabel("$\\psi(x)$")
axes[1].set_xlabel("$p$")
axes[1].set_ylabel("$\\phi(p)$")

# initialization of plots:
def init():
    la.set_data([], [])
    lb.set_data([], [])
    return la, lb

# update plots once per frame:
def animate(i):
    s = sig(i)
    la.set_data(xf, psi(xf, s))
    lb.set_data(xf, phi(xf, s))
    return la, lb

#run the animation:
anim = FuncAnimation(
    fig,
    animate,
    frames=200,
    init_func=init,
    interval=frame_ms,
    blit=False,
    repeat=True,
)

plt.tight_layout()
plt.show()
\end{python}
The main features are the initial setup of animation parameters and
setting up the (empty) plots.  The init function is called at the
start of the animation, but does little in this example.  In the
animate function, which is called once per frame, the plots are
updated with the data for that frame.  Last, the animation is called
via FuncAnimation.

\section{Homework Problems}

\noindent
{\bf Problem 1:} Get the example program (uncertainty.ipynb) running
on your system successfully.  You have explicit permission to use the
internet, AI, knowledgeable colleagues, etc, in order to accomplish
this.  Please submit your working code (even if unchanged) for this
problem, and include some brief comments about the installation steps
that were needed.  Hopefully, the example code will just work out of
the box for you, in which case this is a very easy problem!  \\[5pt]

\noindent
{\bf Problem 2:} Animate the time evolution of a stationary wave
packet ($\beta=0$) Show the magnitude as a black line, and real and imaginary
components as blue and red lines.  Use the parameterization from our derivation!\\[5pt]

\noindent
{\bf Problem 3:} Animate the time evolution of a travelling wave
packet.  Use the parameterization from our derivation.\\[5pt]

\end{document}
